\chapter{Methodology}

\section{Introduction to the Methodology}
% Write 1–2 paragraphs that:
% - Explain why this chapter exists (to describe how the research was carried out).
% - Summarize the main components (research design, tools, implementation, validation).
% - Link back to the Problem Statement and Literature Review: this chapter shows *how* you addressed the problem.

This chapter explains how the research was carried out and how the 
core algorithms were developed. The synthesis did not arise from a 
gradual accumulation of partial solutions, but from a decisive 
insight gained after years of experience with \(3\)D illustration. 
That insight was to frame the entire problem in terms of affine 
linear algebra, which provided the clarity and rigor needed to 
resolve clipping and occlusion systematically.

The implementation was carried out in \LaTeX{} and Lua, not by 
accident but because I am a mathematics textbook illustrator, and 
these are the tools I know and use daily. This choice ensured that 
the algorithms were embedded directly in the environment where they 
are most relevant, while also demonstrating their applicability 
beyond it. The following sections present the design, implementation, 
and validation of the algorithms, showing how they directly address 
the fundamental gaps identified in the Problem Statement and 
Literature Review.

\section{Research Design and Overall Approach}
% Write 1–2 paragraphs that:
% - State the overall research design (qualitative, quantitative, mixed-methods, computational, experimental).
% - Justify why this design fits your research objectives.
% - Mention whether the approach is exploratory, confirmatory, or developmental.

\section{Detailed Description of the Proposed Approach}
% Write several paragraphs that:
% - Describe the methodology, algorithm, framework, or system step by step.
% - Break it into subcomponents if necessary (e.g., preprocessing, computation, output).
% - Keep the description precise but readable: focus on the "what" before the "why."

\section{Rationale and Development Process of the Approach}
% Write 1–2 paragraphs that:
% - Explain the reasoning behind choosing this method (why this over alternatives).
% - Discuss influences from existing literature or practice.
% - If you adapted or extended something, explain how and why.

\section{Data, Tools, and Resources Used}
% Write 1–2 paragraphs that:
% - Identify datasets (source, size, characteristics).
% - List software libraries, frameworks, or languages.
% - Note hardware or computational resources if relevant.
% - Explain why these resources were chosen (availability, reliability, suitability).

\section{Implementation Details}
% Write several paragraphs (with equations, pseudocode, or diagrams if appropriate) that:
% - Provide enough technical detail for reproducibility.
% - Specify parameters, configurations, or assumptions used.
% - Use structured elements (tables, figures, algorithms) for clarity where needed.

\section{Validation Strategy}
% Write 1–2 paragraphs that:
% - Describe how you plan to test or evaluate the approach.
% - Define benchmarks, comparison methods, or criteria for success.
% - Explain whether validation is internal (self-testing), external (against existing methods), or both.

\section{Challenges Encountered and Solutions Implemented}
% Write 2–3 paragraphs that:
% - Describe obstacles (technical, theoretical, practical).
% - Explain how you addressed them (workarounds, adjustments, innovations).
% - Reflect on what these solutions reveal about the strengths/weaknesses of the method.

\section{Limitations of the Current Methodology}
% Write 1–2 paragraphs that:
% - Acknowledge shortcomings, constraints, or simplifying assumptions.
% - Be transparent but constructive: limitations set up justification for future work.

\section{Remaining Challenges and Directions for Future Work}
% Write 1–2 paragraphs that:
% - Point out unresolved methodological questions.
% - Suggest how future research could refine or extend the approach.
% - Distinguish between what you couldn’t address due to scope vs. what remains genuinely open.

\section{Summary of the Methodology}
% Write 1 short paragraph that:
% - Recaps the key methodological choices.
% - Restates why this methodology is appropriate for the research problem.
% - Provides a smooth transition into the Results and Analysis chapter.
